% Part: lambda-calculus
% Chapter: introduction
% Section: reduction

\documentclass[../../../include/open-logic-section]{subfiles}

\begin{document}

\olfileid{lam}{int}{red}
\olsection{लॅम्डा पदांचे न्यूनीकरण}

लॅम्डा पदांचे काय करता येते? ती सोपी करता येतात. $M$ आणि $N$ ही
कोणतीही लॅम्डा पदे आणि $x$ हे कोणतेही चर असो. आदेशनानंतर $N$ मधील
मुक्त चरांशी अडथळा निर्माण करतील अशा $M$ मधील बद्ध चरांची आधी
नावे बदलून, मग $M$ मध्ये~$x$ च्या जागी $N$ आदेशल्यास जो परिणाम
मिळतो, तो आपण $\Subst{M}{N}{x}$ ने दर्शवू. उदाहरणार्थ,
\[
\Subst{(\lambd[w][xxw])}{yyz}{x} = \lambd[w][(yyz)(yyz)w].
\]

\begin{digress}
आदेशनासाठी $[N/x]M$, $[x/N]M$ आणि $M[x/N]$ ही पर्यायी
संकेतनेही वापरली जातात. गोंधळ टाळा!
\end{digress}

सहजप्रज्ञेने पाहता, $(\lambd[x][M])N$ आणि $\Subst{M}{N}{x}$
यांचा अर्थ एकच आहे; पहिले पद बदलून दुसरे करण्याच्या क्रियेला
\emph{$\beta$-संकुचन} म्हणतात. $(\lambd[x][M])N$ ला
\emph{रेडेक्स}, आणि $\Subst{M}{N}{x}$ ला त्याचे
\emph{संकुचनफल} म्हणतात. सर्वसाधारणपणे, एखाद्या उपपदाचे
$\beta$-संकुचन करून $P$ हे पद~$P'$ असे बदलता येत असेल, तर
$P$ चे \emph{एका पायरीत $\beta$-न्यूनीकरण होऊन $P'$ मिळते}
असे म्हणतो आणि $P \redone P'$ लिहितो. $P$ पासून अशा
एक-पायरी न्यूनीकरणांच्या काही पायऱ्यांनी (एकही नसेल तरी)
~$P'$ मिळत असेल, तर $P$ चे~$P'$ पर्यंत
\emph{$\beta$-न्यूनीकरण होते}; संकेतांत, $P \red P'$.
ज्या पदाचे आणखी $\beta$-न्यूनीकरण करता येत नाही, त्याला
\emph{$\beta$-अन्यूनीकरणीय} किंवा \emph{$\beta$-सामान्य}
म्हणतात. संदर्भ स्पष्ट असेल तेव्हा “$\beta$-न्यूनीकरण होते”
इत्यादींऐवजी फक्त “न्यूनीकरण होते” असे म्हणू.

काही उदाहरणे पाहू.
\begin{enumerate}
\item आपल्याला मिळते:
\begin{align*}
(\lambd[x][xxy]) \lambd[z][z] & \redone (\lambd[z][z])(\lambd[z][z]) y \\
& \redone (\lambd[z][z]) y \\
& \redone y.
\end{align*}
\item पद “सोपे” केल्याने ते अधिक गुंतागुंतीचे होऊ शकते:
\begin{align*}
(\lambd[x][xxy])(\lambd[x][xxy]) & \redone (\lambd[x][xxy])(\lambd[x][xxy])y \\
& \redone (\lambd[x][xxy])(\lambd[x][xxy])yy \\
& \redone \dots
\end{align*}
\item पद तसेच राहूही शकते:
\[
(\lambd[x][xx])(\lambd[x][xx]) \redone (\lambd[x][xx])(\lambd[x][xx]).
\]
\item शिवाय, काही पदांचे एकाहून अधिक प्रकारे न्यूनीकरण करता येते;
  उदाहरणार्थ,
\[
(\lambd[x][(\lambd[y][yx]) z)] v \redone (\lambd[y][yv]) z
\]
हे सर्वांत बाहेरच्या उपयोजनाचे संकुचन केल्याने मिळते; आणि
\[
(\lambd[x][(\lambd[y][yx]) z)] v \redone (\lambd[x][zx]) v
\]
हे सर्वांत आतल्या उपयोजनाचे संकुचन केल्याने मिळते. मात्र येथे
लक्षात घ्या की दोन्ही पदांचे पुढे न्यूनीकरण होऊन तेच पद, $zv$, मिळते.
\end{enumerate}

शेवटच्या उदाहरणात मिळणारा एकच अंतिम परिणाम हा योगायोग नाही;
तो लॅम्डा कलनाचा “चर्च--रॉसर गुणधर्म” म्हणून ओळखला जाणारा
सखोल आणि महत्त्वाचा गुणधर्म दाखवतो.

\end{document}
